\subsection{Bid-Ask Spread Decomposition}

\begin{definition} \hlt{Spread Notation}\\
Let $b_t$ be the best bid, $a_t$ the best ask, and $m_t=(a_t+b_t)/2$ the quoted midpoint. The quoted spread is $s_t=a_t-b_t$, the half-spread is $h_t=s_t/2$, and the relative spread is $s_t/m_t$.\\
Cross-sectional spread work should usually compare relative spreads, not dollar spreads.
\end{definition}

\begin{taxonomy} \hlt{Spread Measurement Segmentation}
\begin{enumerate}[label=\roman*.]
\setlength{\itemsep}{0pt}
\item \hlt{quoted spread}: posted ask minus posted bid;
\item \hlt{inside spread}: best ask across dealers or venues minus best bid;
\item \hlt{outside spread}: value trader's sell threshold minus buy threshold;
\item \hlt{effective spread}: execution price relative to the contemporaneous midpoint;
\item \hlt{realised spread}: execution price relative to the midpoint after later price adjustment.
\end{enumerate}
Quoted spreads measure offered liquidity.\\
Realised spreads measure what liquidity suppliers keep after prices move.
\end{taxonomy}

\begin{definition} \hlt{Outside Spread}\\
If a value trader estimates $\widehat{V}$, buys only below $b^o$, and sells only above $a^o$, then $s^o=a^o-b^o$ is the \hlt{outside spread}. It is a private willingness interval, not necessarily a displayed quote. It widens with valuation uncertainty, research cost, impact, financing cost, and adverse-selection risk.
\end{definition}

\begin{taxonomy} \hlt{Spread Component Segmentation}
\begin{enumerate}[label=\roman*.]
\setlength{\itemsep}{0pt}
\item \hlt{order-processing}: systems, clearing, capital, and operating costs;
\item \hlt{inventory}: compensation for warehousing price and funding risk;
\item \hlt{adverse selection}: compensation for trading before prices adjust to informed flow.
\end{enumerate}
\end{taxonomy}

\begin{model} \hlt{Simple Dealer Quote Decomposition}\\
Let $V_0$ be the dealer's current value estimate. Let $V^B$ be the value estimate conditional on the next trader being a buyer, and $V^S$ the estimate conditional on the next trader being a seller. If $C$ is the transaction-cost component and $A=V^B-V^S$ is the adverse-selection component, then $a=V^B+\frac{C}{2}$, $b=V^S-\frac{C}{2}$, $a-b=A+C$.\\
The transaction-cost component is the part a dealer expects to keep in uninformed two-sided flow. The adverse-selection component is the price revision forced by the information in trade direction.
\end{model}

\begin{proposition} \hlt{Adverse-Selection Component Equivalence}\\
Suppose $V_0=(V^B+V^S)/2$ and $V^B\geq V^S$. The adverse-selection spread component $A=V^B-V^S$ is both the gap between conditional value estimates and twice the dealer's mark-to-market loss from quoting around $V_0$ against a correctly signed informed trade.
\begin{proof}
The information view gives $A=V^B-V^S$. If the next buyer is informed and the dealer sells around $V_0$, the value revision is $V^B-V_0=A/2$. If the next seller is informed, the symmetric revision is $V_0-V^S=A/2$. Thus the accounting loss per side is one half of the same component.
\end{proof}
\end{proposition}

\begin{principle} \hlt{Quoted vs Realised Spread}\\
The quoted spread is an offer. The realised spread is smaller when prices move against the liquidity supplier after the trade, which is the empirical footprint of inventory pressure or adverse selection.
\end{principle}

\begin{model} \hlt{Effective Spread Decomposition}\\
For trade sign $\epsilon_t\in\{+1,-1\}$, execution price $p_t$, midpoint $m_t$, and later midpoint $m_{t+\Delta}$, $ES_t=2\epsilon_t(p_t-m_t)$, $RS_t=2\epsilon_t(p_t-m_{t+\Delta})$, $PI_t=2\epsilon_t(m_{t+\Delta}-m_t)$. Then $ES_t=RS_t+PI_t$. $RS_t$ estimates what the liquidity supplier keeps after price adjustment; $PI_t$ estimates the signed price revision associated with the trade.
\end{model}

\begin{mechanism} \hlt{Component Identification}\\
Transaction-cost spread creates bid-ask bounce: trade prices reverse between bid and ask when value is unchanged. Adverse selection creates signed midpoint revision: buyer-initiated trades are followed by higher midpoints and seller-initiated trades by lower midpoints.\\
Empirical decompositions therefore require trade signs and future quote revisions, not prints alone.
\end{mechanism}

\begin{proposition} \hlt{Roll Serial-Covariance Spread Estimator}\\
Assume transaction price $P_t=V_t+\frac{S}{2}Q_t$, where $V_t=V_{t-1}+\epsilon_t$, $\mathbb{E}[\epsilon_t]=0$, $\operatorname{Var}(\epsilon_t)=\sigma^2$, and $Q_t\in\{+1,-1\}$ is an independent buyer/seller sign with equal probabilities. Then $\operatorname{Var}(\Delta P_t)=\sigma^2+\frac{S^2}{2}$, $\operatorname{Cov}(\Delta P_t,\Delta P_{t-1})=-\frac{S^2}{4}$.\\
Hence $S=2\sqrt{-\operatorname{Cov}(\Delta P_t,\Delta P_{t-1})}$ when the serial covariance is negative.
\begin{proof}
Since $\Delta P_t=\epsilon_t+\frac{S}{2}(Q_t-Q_{t-1})$, independence gives $\operatorname{Var}(\Delta P_t)=\sigma^2+\frac{S^2}{4}\mathbb{E}[(Q_t-Q_{t-1})^2]=\sigma^2+\frac{S^2}{2}$.\\
The innovation terms have zero cross-covariance, so $\operatorname{Cov}(\Delta P_t,\Delta P_{t-1}) =\frac{S^2}{4}\mathbb{E}[(Q_t-Q_{t-1})(Q_{t-1}-Q_{t-2})] =-\frac{S^2}{4}$.\\
Solving gives the estimator.
\end{proof}
\end{proposition}

\begin{caveat} \hlt{Roll Estimator Limits}\\
Roll's estimator identifies spread-induced transitory volatility only under strong assumptions. Order splitting, inventory control, serially correlated trade signs, price discreteness, hidden liquidity, and informed trading can make the covariance less negative or positive.
\end{caveat}

\begin{principle} \hlt{Uninformed Trading Loss Channels}\\
Uninformed market orders pay for immediacy through the spread. Uninformed limit orders save the spread only when they execute, and execution is most likely when the quote has become attractive to better-informed flow. The order type changes the loss channel; it does not remove adverse selection.
\end{principle}

\begin{taxonomy} \hlt{Primary Spread Determinants}
\begin{enumerate}[label=\roman*.]
\setlength{\itemsep}{0pt}
\item \hlt{asymmetric information}: raises the adverse-selection component;
\item \hlt{volatility}: raises inventory risk, stale-quote risk, and the value of timing options;
\item \hlt{utilitarian trading interest}: lowers spreads when activity lets dealers and limit-order traders spread fixed costs over many trades.
\end{enumerate}
\end{taxonomy}

\begin{taxonomy} \hlt{Outside Spread Determinant Segmentation}
\begin{enumerate}[label=\roman*.]
\setlength{\itemsep}{0pt}
\item \hlt{valuation error}: wider when fair value is hard to estimate or model errors are correlated;
\item \hlt{winner's curse}: wider when execution itself reveals the trader may be too optimistic or pessimistic;
\item \hlt{research recovery}: wider when few analysed instruments become tradable opportunities;
\item \hlt{position and financing}: wider when expected holding period, leverage, borrow, or funding costs are large;
\item \hlt{market impact}: wider when entering or exiting the position moves price.
\end{enumerate}
Outside spreads are usually wider than dealer spreads because value traders hold risk longer and cannot recycle capital through balanced two-sided flow.
\end{taxonomy}

\begin{taxonomy} \hlt{Order-Driven Spread Determinant Segmentation}
\begin{enumerate}[label=\roman*.]
\setlength{\itemsep}{0pt}
\item \hlt{information asymmetry}: large effect, usually widening spreads;
\item \hlt{cancellation speed}: slow repricing widens spreads by increasing stale-quote risk;
\item \hlt{volatility}: widens spreads through both risk and timing-option value;
\item \hlt{limit-order management cost}: widens spreads when passive liquidity is expensive to monitor;
\item \hlt{trader time value}: widens spreads when waiting is costly;
\item \hlt{commission difference}: passes through roughly one-for-one into the spread;
\item \hlt{risk aversion}: widens spreads when passive execution risk is disliked.
\end{enumerate}
\end{taxonomy}

\begin{taxonomy} \hlt{Secondary Spread Proxy Segmentation}
\begin{enumerate}[label=\roman*.]
\setlength{\itemsep}{0pt}
\item \hlt{information proxies}: disclosure quality, analyst coverage, information vendors, firm age, business complexity, insider-trading enforcement, and expected material news;
\item \hlt{volatility proxies}: leverage, operating uncertainty, storage or inventory risk, macro rate risk, credit risk, and growth uncertainty;
\item \hlt{utilitarian-interest proxies}: volume, trade count, firm size, debt issue size, hedging demand, and availability of close risk substitutes.
\end{enumerate}
These are proxies, not primitives.\\
A variable can load on more than one determinant, so cross-sectional tests need controls.
\end{taxonomy}

\begin{caveat} \hlt{Liquidity and Capital Structure}\\
When asymmetric information is severe and natural trading interest is weak, the spread required for public trading can be too wide for a useful market.\\
Private debt, venture capital, relationship lending, and monitored equity can dominate public trading because they replace anonymous liquidity with due diligence and control rights.
\end{caveat}
